% Generated draft: requires mathematical and source review.
\documentclass[11pt]{article}
\usepackage[margin=25mm]{geometry}
\usepackage{fontspec}
\setmainfont{Latin Modern Roman}
\usepackage{amsmath,amssymb,mathtools,unicode-math}
\setmathfont{Latin Modern Math}
\usepackage{enumitem,needspace,xurl,hyperref}
\setlength{\parindent}{0pt}
\setlength{\parskip}{5pt}
\setlength{\emergencystretch}{3em}
\newcommand{\N}{\mathbb{N}}
\newcommand{\Z}{\mathbb{Z}}
\newcommand{\Q}{\mathbb{Q}}
\newcommand{\R}{\mathbb{R}}
\newcommand{\C}{\mathbb{C}}
\newcommand{\F}{\mathbb{F}}
\newcommand{\A}{\mathbb{A}}
\newcommand{\PP}{\mathbb{P}}
\newcommand{\E}{\mathbb{E}}
\newcommand{\eps}{\varepsilon}
\newcommand{\dd}{\,\mathrm{d}}
\begin{document}
\section*{UnsolvedMath ID: 30001448}
\textbf{Source problem code:} catalogue entry
\subsection{Definitions and mathematical statement}
In standard Coxeter terminology, for a Coxeter system \((W,S)\) and \(J\subseteq S\), the standard parabolic subgroup is
\[
W_J=\langle s:s\in J\rangle.
\]
A maximal proper standard parabolic corresponds to \(J=S\setminus\{s\}\) for some \(s\in S\); a non-maximal proper standard parabolic has \(J\subsetneq S\) that is not maximal among proper subsets. These conventions explain possible terminology only: the supplied record does not specify the intended Coxeter system, parabolic modules, or meaning of ``maximally parabolic.'' Kazhdan--Lusztig elements are the distinguished elements arising from the relevant Kazhdan--Lusztig construction; the construction and normalization intended here are not supplied.

The source poses two research questions:
\begin{enumerate}
\item Find a combinatorial interpretation for evaluations of Kazhdan--Lusztig elements for non-maximal parabolic modules.
\item Find a Cauchy summation formula for the Kazhdan--Lusztig elements of a maximally parabolic module.
\end{enumerate}
The evaluations in the first question are unspecified. The summation indices, variables, coefficients, pairing of elements, and proposed right-hand side in the second question are also unspecified. Consequently, these are source-level research requests, not a uniquely determined quantified conjecture or identity. In particular, no assertion over all Coxeter systems or all parabolic modules is implied by this draft.
\subsection{Short English statement}
The problem asks for combinatorial interpretations of evaluations of non-maximal parabolic Kazhdan--Lusztig elements and for a Cauchy summation identity in a maximally parabolic setting. The supplied record omits the algebraic setting and the evaluations and identity to be studied, so further source context is needed for a fully precise formulation.
\subsection{Sources}
\begin{itemize}
\item Supplied local UnsolvedMath registry record, permanent ID \(30001448\), title ``Combinatorics and Cauchy Identities for Parabolic KL Elements,'' source problem code \texttt{catalogue entry}, in the Oberwolfach Reports Open Problems collection.
\item Supplied authoritative source citation: \emph{Combinatorial Representation Theory} (2010), Oberwolfach Reports, pp.\ 799--882, DOI \texttt{10.4171/owr/2010/15}. Source URL and supplied literature reference labelled ``original OWR report'': \url{https://doi.org/10.4171/owr/2010/15}. The registry gives the proposed year as 2010 and supplies no named proposer. Its background describes Oberwolfach Reports as documenting mathematical workshops and their open-problem discussions.
\item The registry labels the problem \texttt{open} and the statement \texttt{corrected\_verified}. Its supplied verification note says that the mathematical content was checked against the official report and that notation, context, and question form were normalized without changing the source-level claim. These are supplied curation assertions, not independent verification in this draft.
\item Literature assessment dated exactly \texttt{2026-08-21}: the supplied curation states that the record lists two open Kazhdan--Lusztig-element problems, concerning non-maximal parabolic combinatorics and a maximally parabolic Cauchy identity, and that no complete general solution was found. The listed source checked is the original OWR report linked above. This is supplied dated literature triage, not an independently verified assessment; the record explicitly requests independent verification of time-sensitive status and cited claims.
\end{itemize}
\end{document}
